\documentclass[conference]{IEEEtran}
\IEEEoverridecommandlockouts
% Packages: only those installed in the local MiKTeX and actually used (checked with kpsewhich).
% Removed from the course template because they are missing locally and unused here:
% algorithmic, algorithm2e, subfigure, natbib, multirow, threeparttable, diagbox.
\usepackage{cite}
\usepackage{amsmath,amssymb,amsfonts}
\usepackage{graphicx}
\usepackage{textcomp}
\usepackage{xcolor}
\usepackage{bm}
\usepackage{booktabs}
\usepackage{url}

% ---------------------------------------------------------------- notation
\newcommand{\vh}{\bm{h}}
\newcommand{\hh}{\hat{\bm{h}}}
\newcommand{\vy}{\bm{y}}
\newcommand{\vn}{\bm{n}}
\newcommand{\vmu}{\bm{\mu}}
\newcommand{\mC}{\bm{C}}
\newcommand{\mI}{\bm{I}}
\newcommand{\cP}{\mathcal{P}}
\newcommand{\cM}{\mathcal{M}}
\newcommand{\CN}{\mathcal{CN}}
% Draft markers: every [TBD] must be gone before submission (grep for "TBD").
\newcommand{\tbd}[1]{\textbf{[TBD: #1]}}
% Figure slot: includes figures/<file> once it exists, otherwise draws a framed placeholder of
% the planned size, so the page count of the draft stays realistic.
\newcommand{\figslot}[4]{% #1 file, #2 width, #3 height, #4 placeholder text
  \IfFileExists{figures/#1}{\includegraphics[width=#2]{figures/#1}}{%
    \fbox{\parbox[c][#3][c]{\dimexpr#2-2\fboxsep-2\fboxrule\relax}{\centering\footnotesize
      \textbf{[TBD: \texttt{\detokenize{figures/#1}}]}\\[2pt] #4}}}}

\begin{document}

\title{OFDM Channel Reconstruction from Sparse Noisy Pilots Using Measured CSI}

% Names: pinyin of the names in docs/proposal/build/build_proposal.py (MEMBERS) -- to be confirmed.
\author{\IEEEauthorblockN{Zhengyang Feng (225010324), Hao Xie (226010063), Hengyue San (226010043),
Peiyu Zhong (225010390)}
\IEEEauthorblockA{Group 12, School of Science and Engineering, The Chinese University of Hong Kong, Shenzhen\\
\small{E-mails: [TBD]}}}

\maketitle

\begin{abstract}
A pilot-based OFDM receiver must infer the channel on subcarriers without pilots, and its accuracy limits
how far the pilot overhead can be reduced. Learned estimators are often reported to beat classical ones,
mostly on simulated channels. We reconstruct measured channels of the public DICHASUS dichasus-005x
dataset (32 receive antennas, indoor non-line-of-sight) from a few noisy pilots. On a 64-subcarrier block
we compare linear and spline interpolation, LMMSE with training-data statistics, and multilayer
perceptrons (MLPs) with and without knowledge of the noise level, for 0--30\,dB pilot SNR and 16, 8 or 4
pilots; for joint antenna--frequency reconstruction we compare a Gaussian-mixture estimator and a joint
MLP with joint LMMSE. Selection and reporting use separate blocks of a validation run with paired
block-bootstrap confidence intervals, and a held-out run is evaluated once.
On the held-out run no MLP variant is better than LMMSE in any of the 84 settings of the 64-subcarrier
block: all are within about 0.2\,dB up to 15\,dB, and the plain MLP is 0.9--1.6\,dB worse at 30\,dB. In the
joint setting the Gaussian-mixture estimator is better than joint LMMSE in all eight settings, by
0.2--2.1\,dB with pilots on all antennas and by 2.9--5.9\,dB with pilots on 8 of 32 antennas, at up to 64
times the multiply--accumulate count.
\end{abstract}

\begin{IEEEkeywords}
OFDM channel estimation, pilot interpolation, LMMSE, measured CSI, neural networks, Gaussian mixture model
\end{IEEEkeywords}

%=====================================================================================================
\section{Introduction}
An OFDM receiver estimates the channel at a few pilot subcarriers and must infer it on all others. The
accuracy is limited by the noise on the pilot estimates and by the pilot spacing, while every pilot
occupies a resource element that could carry data. Classical receivers interpolate least-squares pilot
estimates~\cite{coleri2002} or apply linear minimum mean-square-error (LMMSE) estimation with the
frequency correlation of the channel~\cite{edfors1998}. Neural estimators such as
ChannelNet~\cite{soltani2019} and ReEsNet~\cite{li2020} reported gains over these methods, but were mostly
trained and evaluated on simulated channels, and how fair such comparisons are depends strongly on the
baseline and the protocol~\cite{feriani2023}. Measured channels, such as those recorded by the DICHASUS
channel sounder~\cite{euchner2021sounder,euchner2021data}, allow a more realistic test.

We ask: \emph{under which conditions does a learned estimator help over LMMSE for reconstructing OFDM
channels from noisy sparse pilots?} We vary the pilot SNR, the pilot density, whether the receiver knows
the noise level, and the dimension of the reconstruction (frequency only, or antennas and frequency
jointly). Our contributions are:
\begin{itemize}
\item a leakage-free protocol on measured CSI (one global scale, screening from the CSI itself, split by
run, separate selection and reporting blocks, paired block bootstrap, one test evaluation);
\item a comparison of interpolation, LMMSE and four MLP variants over 0--30\,dB and 16/8/4 of 64 pilots
with three seeds, including noise-power mismatch and computational cost;
\item an extension to joint reconstruction over antennas and subcarriers: joint LMMSE, a Gaussian-mixture
(GMM) estimator~\cite{turan2022} and a joint MLP.
\end{itemize}
The code and every result record are public at \url{https://github.com/zzoewlxx/mce5901-ofdm-channel-reconstruction}.

%=====================================================================================================
\section{Related Work}
\emph{Classical estimation.} Coleri et al.~\cite{coleri2002} compared pilot arrangements and the
interpolation of least-squares pilot estimates. Edfors et al.~\cite{edfors1998} derived the LMMSE
estimator from the frequency correlation of the channel and a low-rank approximation of it by singular
value decomposition. LMMSE is the best \emph{linear} estimator given the first two moments; it coincides
with the MMSE estimator when the channel and the noise are jointly Gaussian, but not in general.

\emph{Learned estimators.} ChannelNet~\cite{soltani2019} treats the pilot grid as a low-resolution image
and applies super-resolution and denoising networks; ReEsNet~\cite{li2020} uses a residual convolutional
network. Both were evaluated mainly on simulated channel models. CeBed~\cite{feriani2023} benchmarks deep
data-driven OFDM channel estimators, including ReEsNet; its public code computes the LMMSE reference from
each sample's true channel, a genie rather than a practical receiver, whereas we fit all statistics on
training data.

\emph{Statistical learning on measured data.} Turan et al.~\cite{turan2022} fit a Gaussian mixture model
to training channels and combine per-component LMMSE estimates; on measured multi-antenna channels it
outperformed LMMSE with the sample covariance. Baur et al.~\cite{baur2024} trained a variational
autoencoder-based estimator on real measurements, which outperformed LMMSE, a CNN estimator and other
baselines. Both works estimate multi-antenna (spatial) channels; we ask whether such gains exist across
OFDM subcarriers, and across antennas and subcarriers jointly.

%=====================================================================================================
\section{Problem Formulation}
\label{sec:problem}
\emph{Signal model.} Let $\vh\in\mathbb{C}^N$ be the channel of one receive antenna on a block of $N=64$
consecutive subcarriers (in the joint setting of Section~\ref{sec:methods}, the stacked
antenna--subcarrier array, $N=512$). Pilots occupy the index set $\cP$ with $|\cP|=n_p$; the remaining set
$\cM$ is withheld. With unit-power pilots, the least-squares pilot estimates are
\begin{equation}
\vy=\vh_\cP+\vn,\qquad \vn\sim\CN(\bm{0},\sigma^2\mI_{n_p}).
\label{eq:model}
\end{equation}
One constant $c$ scales all channels; it is computed on the training set $\mathcal{T}$ so that the
average power per entry is one there, $\sum_{i\in\mathcal{T}}\|c\,\vh_i\|^2=N|\mathcal{T}|$. The pilot SNR
is then $1/\sigma^2$, i.e., $\sigma^2=10^{-\mathrm{SNR}/10}$. Normalizing each snapshot by its
own power is not allowed: it reveals $\|\vh\|$, which a receiver does not know, and in our pilot check it
inflated the MLP.

\emph{Metric.} For an evaluation set $\mathcal{E}$ we report
\begin{equation}
\mathrm{NMSE}=10\log_{10}\frac{\sum_{i\in\mathcal{E}}\|\hh_{\cM,i}-\vh_{\cM,i}\|^2}
{\sum_{i\in\mathcal{E}}\|\vh_{\cM,i}\|^2}\ \text{dB},
\label{eq:nmse}
\end{equation}
on the withheld entries only, since the pilot entries are observed directly.

\emph{LMMSE.} With the sample mean $\vmu$ and covariance $\mC$ of the training channels,
\begin{equation}
\hh=\vmu+\mC_{:,\cP}\big(\mC_{\cP\cP}+(\sigma^2+\rho\,\bar p)\mI\big)^{-1}(\vy-\vmu_\cP),
\label{eq:lmmse}
\end{equation}
where $\bar p=\mathrm{tr}(\mC_{\cP\cP})/n_p$ and $\rho\ge 0$ is a small diagonal loading. For centered
training channels $\tilde\vh_i=\vh_i-\vmu$, averaging the squared error of a linear map $\bm{W}$ over the
noise gives $\sum_i\|\tilde\vh_i-\bm{W}\tilde\vh_{\cP,i}\|^2+|\mathcal{T}|\sigma^2\|\bm{W}\|_F^2$. This is
ridge regression from the clean pilots to the channel; its minimizer is the matrix in \eqref{eq:lmmse}
with $\rho=0$, and the loading $\rho$ adds shrinkage. The expected error of \eqref{eq:lmmse} on $\cM$ is
$\varepsilon(\sigma^2)=\mathrm{tr}\big(\mC_{\cM\cM}-\mC_{\cM\cP}(\mC_{\cP\cP}+\sigma^2\mI)^{-1}
\mC_{\cP\cM}\big)$ (for $\rho=0$).

\emph{GMM estimator.} Under the mixture prior $p(\vh)=\sum_{k=1}^K\pi_k\,\CN(\vh;\bm{0},\mC_k)$ and the
model \eqref{eq:model}, the conditional mean is
\begin{equation}
\hh=\sum_{k=1}^K p(k\,|\,\vy)\,\mC_{k,:\cP}\big(\mC_{k,\cP\cP}+\sigma^2\mI\big)^{-1}\vy,
\label{eq:gmm}
\end{equation}
with $p(k\,|\,\vy)\propto\pi_k\,\CN(\vy;\bm{0},\mC_{k,\cP\cP}+\sigma^2\mI)$~\cite{turan2022}. For $K=1$
it reduces to \eqref{eq:lmmse} with zero mean.

\emph{Learned estimator.} An MLP $f_\theta:\mathbb{R}^{2n_p+1}\to\mathbb{R}^{2|\cM|}$ maps
$[\mathrm{Re}\,\vy;\mathrm{Im}\,\vy;\mathrm{SNR}/30]$ to $[\mathrm{Re}\,\hh_\cM;\mathrm{Im}\,\hh_\cM]$
and is trained to minimize the squared error. Because the conditional mean
$\mathbb{E}[\vh_\cM\,|\,\vy]$ is in general nonlinear in $\vy$ when $\vh$ is not Gaussian, a learned or
mixture estimator can in principle beat \eqref{eq:lmmse}; whether the measured channel departs enough
from a Gaussian is the empirical question.

\emph{Assumptions.} (A1)~Pilot errors are additive white circular Gaussian noise; the noise is synthetic,
and the measured CSI, which has its own estimation error, serves as the reference. (A2)~The receiver knows
$\sigma^2$ unless stated otherwise (Section~\ref{sec:setup} relaxes this). (A3)~Training and test runs
share the channel statistics (same room). (A4)~One snapshot at a time: no correlation over time is used,
and in the main setting each antenna is estimated separately. (A5)~A fixed pilot pattern.

%=====================================================================================================
\section{Dataset and Preprocessing}
\label{sec:data}
\emph{Dataset.} dichasus-005x~\cite{euchner2021data} was recorded with the DICHASUS distributed massive
MIMO channel sounder~\cite{euchner2021sounder}. A transmitter on a vacuum robot moved through a room in
which a metallic barrier blocks the line of sight. The carrier frequency is 1.272\,GHz and the bandwidth
50\,MHz, with 1024 OFDM subcarriers and 32 co-located receive antennas. The data (five robot runs 0053--0057, 6.9\,GB, CC~BY~4.0) contain
for each snapshot a $32\times1024$ complex CSI matrix and an SNR field.

\emph{Preparation.} (i)~A NumPy-only TFRecord parser reads the files, and the subcarrier axis is
\texttt{fftshift}ed into frequency order. (ii)~The published per-antenna timing and phase offsets of the
sounder, $\tau_a$ and $\phi_a$, are compensated by multiplying subcarrier $k$ of antenna $a$ by
$e^{j(2\pi k\tau_a/1024+\phi_a)}$. (iii)~We keep subcarriers 256--319 (64 subcarriers, 3.1\,MHz), inside
the occupied band and away from the band edges and the DC subcarrier. (iv)~Screening: a few vectors are
pure noise, and the dataset's SNR field does not identify them. For each (snapshot, antenna) vector we
estimate the CSI error from fourth-order differences across frequency over subcarriers 128--893 without
the DC region, $\hat\varepsilon=\overline{|\Delta^4H_k|^2}/(70\,\overline{|H_k|^2})$, where 70 is the noise gain of
$\Delta^4$ for white errors and the smooth channel nearly cancels. Vectors with $\hat\varepsilon>-30$\,dB
are removed: 0.44\,\% of the training vectors (runs 0053--0055) and 0.40\,\% of the validation vectors. (v)~The split is by run
(Table~\ref{tab:data}), so adjacent snapshots never fall into different sets. Each antenna vector of a
snapshot is one example. The scale $c$, all statistics and all model weights are fitted on the training
runs only.

\begin{table}[t]
\caption{Data split; vectors counted after screening (validation: even / odd blocks)}
\label{tab:data}
\centering\footnotesize
\setlength{\tabcolsep}{3.5pt}
\begin{tabular}{@{}llrrr@{}}
\toprule
Use & Runs & Snapshots & 64-sc vectors & Joint samples\\
\midrule
Training & 0053--0055 & 15\,940 & 507\,853 & 15\,404\\
Validation & 0056 & 5\,220 & 83\,502 / 82\,871 & 2\,530 / 2\,488\\
Test (used once) & 0057 & 5\,001 & 159\,400 & 4\,807\\
\bottomrule
\end{tabular}
\end{table}

%=====================================================================================================
\section{Methods}
\label{sec:methods}
\emph{Interpolation.} Linear and natural cubic-spline interpolation of $\vy$ need neither statistics nor
the noise level and represent common receiver practice~\cite{coleri2002}. Both are fixed linear maps
$\hh=\bm{W}\vy$.

\emph{LMMSE} \eqref{eq:lmmse} uses $\vmu$ and $\mC$ from all training vectors and the true $\sigma^2$. It is
the reference for every gain we report, because it is the best linear estimator and, with about
$5\times10^5$ training vectors for a $64\times64$ covariance, its statistics are accurate. The loading
$\rho\in\{10^{-6},\dots,10^{-1}\}$ is chosen per pilot density and SNR on the selection blocks.

\emph{MLP variants.} All use two hidden layers of 128 ReLU units (about 33\,000 parameters), one model per
pilot density, trained on the training runs with pilot SNR drawn from $\mathcal{U}[0,30]$\,dB and fresh
noise in every mini-batch (batch 256, Adam, 15 epochs, learning rate $10^{-3}$ decayed to $3\cdot10^{-4}$
and $10^{-4}$); the epoch with the best selection-block NMSE is kept. (i)~\emph{Plain}: squared error, as
in the proposal. (ii)~\emph{Relative-error loss}: each sample's squared error is divided by
$\varepsilon(\sigma_i^2)$, the LMMSE error expected at its SNR, so that high-SNR samples, whose errors are
far smaller, are not drowned out. (iii)~\emph{Residual}: the network outputs a correction that is added to
the LMMSE estimate, so it starts from the strong linear solution. (iv)~\emph{Blind}: the SNR input is
fixed to zero; it is needed for a fair noise-mismatch study. A $3\times256$ network, tried on the validation run only, was worse at 30\,dB and was not carried forward. The input is a short vector, not a time--frequency image, so we use fully
connected networks rather than the CNNs of~\cite{soltani2019,li2020}.

\emph{Joint antenna--frequency extension.} Within the 64-subcarrier block the training covariance is
concentrated in very few eigenmodes (the first carries 99.5\,\% of the energy, the first two
99.99\,\%), leaving little room beyond \eqref{eq:lmmse}. Across the 32 antennas many more eigenmodes
matter (17, 28 and 43 of 512 for 90, 99 and 99.9\,\% of the joint energy), yet the channel is governed by few physical quantities (transmitter position and orientation): it
lies near a low-dimensional nonlinear set that one global covariance cannot represent, but a mixture of
local covariances or a network can. We stack all 32 antennas on 16 of the 64 subcarriers (about every
fourth) into $N=512$ entries per snapshot (snapshots whose 32 vectors all pass the screening) and place
pilots on 32 antennas$\times$4 subcarriers or on 8 antennas (every fourth)$\times$4 subcarriers. Methods:
\emph{joint LMMSE} \eqref{eq:lmmse} with the $512\times512$ training covariance; the \emph{GMM estimator}
\eqref{eq:gmm}, fitted by EM on clean training channels from a $k$-means start, with zero-mean components
(each snapshot has a random common phase) and shrinkage $\mC_k\leftarrow(1-\alpha)\mC_k+\alpha\mC$,
$(K,\alpha)$ chosen on the selection blocks among six pairs with $K$ from 8 to 64 and
$\alpha\in\{0.01,0.05\}$, for each of three $k$-means seeds, with EM run until the relative change of the mean
log-likelihood is below $10^{-5}$ (at most 50 iterations); a \emph{joint MLP} and a joint residual
variant, trained as above with a random common phase per snapshot as augmentation
(two hidden layers of 512 ReLU units, about $7.9\times10^5$ parameters, 150 epochs; the residual
variant's output layer starts at zero, so training starts from joint LMMSE). The joint LMMSE loading is
chosen from $\rho\in\{0,10^{-4},10^{-3},10^{-2}\}$ on the selection blocks.

%=====================================================================================================
\section{Experimental Setup}
\label{sec:setup}
\emph{Operating conditions.} Pilot SNR 0, 5, \dots, 30\,dB (joint setting 0, 10, 20, 30\,dB); 16, 8 or 4
of 64 pilots at $\mathrm{round}(\mathrm{linspace}(0,63,n_p))$, so both block ends are pilots and the
withheld subcarriers are interpolated, not extrapolated.

\emph{Validation protocol.} Run 0056 is cut into blocks of 50 consecutive snapshots. \emph{Even} blocks
are used for every selection (loading $\rho$, $(K,\alpha)$, model configuration, early stopping);
\emph{odd} blocks (52 blocks) for every reported validation number. Because consecutive snapshots are
correlated, confidence intervals come from a paired block bootstrap: the reported blocks are resampled
with replacement ($B\ge1000$ times), the NMSE \eqref{eq:nmse} of both methods is computed on the same resampled
blocks, and the 2.5 and 97.5\,\% percentiles of their dB difference form the 95\,\% interval.

\emph{Fairness.} All methods see identical noise realizations: for a given pilot density and SNR the
noise is drawn with seed $1000n_p+\mathrm{SNR}$ (64-subcarrier setting) or $7919|\cP|+\mathrm{SNR}$ (joint
setting). Every MLP is trained with three seeds and the GMM EM is run with three seeds; we report the seed
average (per-block errors averaged over seeds before forming the NMSE) and the per-seed range.

\emph{Noise-power mismatch.} The receiver assumes an SNR $a$ while the true SNR is $t$, both on
$\{0,5,\dots,30\}$\,dB. LMMSE (with the loading chosen for $a$), the SNR-aware MLPs and the LMMSE part of
the residual variant all receive the same wrong value $a$; every design stays frozen and only the true
noise changes. The blind MLP needs no SNR and linear interpolation is unaffected. We report each method's NMSE
loss relative to its own matched design and relative to matched LMMSE.

\emph{Complexity.} Per estimate, interpolation and LMMSE cost $N n_p$ complex multiply--accumulates
(MACs) as one precomputed matrix product, the MLP $\sum_\ell d_{\ell-1}d_\ell$ real MACs over its layer
widths $d_\ell$, and the GMM with precomputed per-component matrices about $K N n_p$ complex MACs; the
counts as implemented are our primary measure. We also measure the CPU time per vector for calls of
10\,000 vectors on one laptop CPU (Intel Core Ultra~5 125H) with one BLAS thread and the process pinned to
the fastest core.

\emph{Test run.} After all choices are frozen, one dedicated script evaluates run 0057 \emph{once}, with
the training scale and the frozen loadings, mixtures and weights; intervals use all test blocks, and a second
interval over merged 200-snapshot blocks checks their sensitivity to correlation between neighbouring
blocks. Before the test we also froze how results are judged: a learned estimator counts as better than
LMMSE only if its paired interval excludes zero \emph{and} its gain is at least 0.1\,dB. The script runs
only after a rehearsal of the same code has reproduced every validation number, and only once. All numbers
in Section~\ref{sec:results} come from this run. All code is NumPy-based
(scikit-learn only for $k$-means); every run is stored with its configuration and git commit.

%=====================================================================================================
\section{Results and Discussion}
\label{sec:results}
\begin{figure*}[t]
\centering
\figslot{fig_nmse_test.pdf}{\textwidth}{1.9in}{NMSE vs.\ pilot SNR on test run 0057; panels 16 / 8 / 4 of 64
pilots; linear, spline, LMMSE, MLP variants; 95\,\% block-bootstrap CIs.}
\caption{NMSE on the withheld subcarriers versus pilot SNR on the test run 0057, for 16, 8 and 4 of 64
pilots (MLPs: seed averages); the 95\,\% intervals are at most $\pm0.6$\,dB.}
\label{fig:main}
\end{figure*}

\subsection{Main comparison on the test run}
The test run has 101 blocks of at most 50 snapshots (Table~\ref{tab:data}). Fig.~\ref{fig:main} shows that
interpolation is limited by the pilot noise at every SNR: the NMSE of linear interpolation stays 1.8--2.6\,dB
below $-\mathrm{SNR}$. LMMSE is 9.7\,dB better at 0\,dB and 5.9\,dB at 30\,dB with 16 pilots, 7.3 and
3.9\,dB with 8, and 5.0 and 1.7\,dB with 4 (all paired intervals within $\pm0.13$\,dB). It averages the
noise over all pilots, which interpolation between two neighbours cannot do, so its advantage shrinks with
fewer pilots; each halving of the pilots costs LMMSE 2.3--2.7\,dB. The spline is 0.8--1.5\,dB \emph{worse}
than linear interpolation at every SNR: the channel is so smooth within the block that the error is almost
all noise, and the spline weights amplify noise more.

\subsection{Gains of the learned estimators over LMMSE}
Fig.~\ref{fig:gain} shows the paired gains of the four MLP variants over LMMSE. Under the frozen rule
no variant is better than LMMSE in any of the 84 cells (4 variants, 3 densities, 7 SNRs): no interval
lies above zero, and the largest gain is $+0.02$\,dB. 35 cells are worse (interval below zero, loss at
least 0.1\,dB); in the other 49 the whole interval lies within $\pm0.1$\,dB. With 200-snapshot blocks
the same 35 cells are worse and none is better. The plain MLP is within 0.07\,dB of
LMMSE up to 15\,dB and loses 0.12--0.18\,dB at 20\,dB, 0.26--0.50\,dB at 25\,dB and 0.93--1.56\,dB at
30\,dB. The relative-error loss reduces the 30-dB loss to 0.33--0.57\,dB but costs up to 0.22\,dB at low
SNR. The residual MLP equals LMMSE (difference below 0.005\,dB): its correction is trained to almost
zero. Seeds differ by at most 0.36\,dB. This matches Section~\ref{sec:methods}: within the block the
channel is close to a two-dimensional linear family, for which LMMSE with accurate statistics leaves
nothing to learn; the plain MLP's approximation error shows from 20\,dB on, where the LMMSE error is
below about $-25$\,dB.

\begin{figure}[t]
\centering
\figslot{fig_gain_test.pdf}{\columnwidth}{2.0in}{Gain over LMMSE (dB) vs.\ SNR for plain, relative-error,
residual and blind MLP; mean and range over three seeds.}
\caption{Gain of the MLP variants over LMMSE on the test run (dB, positive = better): seed average
with its paired 95\,\% interval.}
\label{fig:gain}
\end{figure}

\subsection{Unknown noise power}
Fig.~\ref{fig:mismatch}(a) shows the loss when the receiver assumes a wrong SNR, averaged over the true
SNR. The loss is asymmetric: assuming an SNR 10\,dB too low costs LMMSE 2.1--2.4\,dB (16, 8, 4 pilots),
10\,dB too high only 0.7\,dB. At the extremes, assuming 0\,dB at a true SNR of 30\,dB costs 16.8--20.0\,dB,
assuming 30\,dB at 0\,dB costs 2.9--3.0\,dB: an SNR assumed too low makes the estimator smooth away
signal, one assumed too high only lets more noise through. The SNR-aware MLPs are only slightly less
sensitive (worst case 13.0--19.3\,dB against their own matched results). In Fig.~\ref{fig:mismatch}(b) an LMMSE designed for 20\,dB loses 2.3\,dB at
0\,dB and 1.3\,dB at 30\,dB. The blind MLP needs no SNR: for all three densities its seed average
stays within 0.29\,dB of the LMMSE that knows the SNR up to 20\,dB, and loses 0.40--0.64\,dB at 25\,dB and
1.1--1.7\,dB at 30\,dB.

\begin{figure}[t]
\centering
\figslot{fig_mismatch_test.pdf}{\columnwidth}{1.9in}{(a) NMSE loss vs.\ SNR offset (assumed $-$ true);
(b) designs for 20\,dB vs.\ true SNR, incl.\ blind MLP; 16 pilots.}
\caption{Wrong noise-power assumption, test run, 16 pilots. (a)~Loss relative to the same method with
the true SNR, averaged over the true SNRs. (b)~LMMSE and MLP given 20\,dB at every true SNR, and the blind MLP: loss
relative to LMMSE with the true SNR.}
\label{fig:mismatch}
\end{figure}

\subsection{Computational cost}
Table~\ref{tab:cost} gives the cost per estimate. The MAC counts are exact; the times come from one run
on the idle laptop (repeats within 25\,\%). With 64
subcarriers every method needs well under 10\,$\mu$s per vector, and time does not follow MACs: the MLP
needs 16 times the MACs of LMMSE but only about 7 times its time. In the joint setting the GMM with $K=64$ needs 64 times the MACs of joint
LMMSE and about 61 times its time, 0.76\,ms per vector, with its per-component matrices precomputed and
stored (8.4 million real parameters).

\begin{table}[t]
\caption{Cost per estimate: 8 of 64 pilots (top), 32$\times$4 joint pilots (bottom)}
\label{tab:cost}
\centering\footnotesize
\setlength{\tabcolsep}{4pt}
\begin{tabular}{@{}lrrr@{}}
\toprule
Method & Real param. & Real MACs & CPU ($\mu$s)\\
\midrule
Linear interpolation & 512 & 2\,048 & 0.29\\
LMMSE & 1\,152 & 2\,048 & 0.57\\
MLP $2\times128$ & 33\,264 & 32\,896 & 3.9\\
LMMSE + residual MLP & 34\,553 & 35\,216 & 5.4\\
\midrule
Joint LMMSE & 132\,096 & 262\,144 & 12.4\\
Joint MLP $2\times512$ & 788\,736 & 786\,944 & 43.5\\
Joint LMMSE + residual MLP & 953\,729 & 1\,114\,880 & 65.8\\
GMM, $K=64$ & 8\,388\,672 & 16\,809\,984 & 763\\
\bottomrule
\end{tabular}
\end{table}

\subsection{Joint antenna--frequency reconstruction}
Table~\ref{tab:joint} gives the joint results. With pilots on all 32 antennas the GMM gains 2.08\,dB
over joint LMMSE at 0\,dB and 0.17\,dB at 30\,dB. With pilots on 8 antennas, joint LMMSE cannot
extrapolate to the 24 antennas without pilots (NMSE between $-2.0$ and $-3.0$\,dB at every SNR), and
the GMM gains 2.95--5.85\,dB. All eight gains are ``better'' under the frozen rule,
also with 200-snapshot blocks, and the three seeds agree within 0.1\,dB. The plain joint MLP gains
2.2--4.7\,dB with $8\times4$ pilots but has an error floor with $32\times4$ pilots (3.6 and 11.1\,dB worse
than LMMSE at 20 and 30\,dB); the residual variant removes the floor. The GMM is better than the plain
joint MLP in all eight rows and better than the residual MLP in seven; in the eighth ($8\times4$,
10\,dB) the difference is inconclusive. Two cautions: with $8\times4$ pilots the error stays high even
with the GMM ($-4.9$ to $-8.8$\,dB), and the selected $K$ is the grid's largest in seven rows.

\begin{table}[t]
\caption{Joint setting, test run: joint LMMSE NMSE (dB), gains over it (dB, seed averages), and the
GMM's paired gain over the residual MLP with 95\,\% interval}
\label{tab:joint}
\centering\footnotesize
\setlength{\tabcolsep}{3pt}
\begin{tabular}{@{}lrrrrrc@{}}
\toprule
Pilots & SNR & LMMSE & GMM & MLP & Res. & GMM $-$ Res.\\
\midrule
$32\times4$ & 0  & $-8.15$ & $+2.08$ & $+0.31$ & $+0.84$ & $+1.24$ [$+1.17$, $+1.31$]\\
            & 10 & $-16.24$ & $+0.78$ & $-0.18$ & $+0.18$ & $+0.60$ [$+0.57$, $+0.63$]\\
            & 20 & $-25.02$ & $+0.39$ & $-3.55$ & $-0.01$ & $+0.40$ [$+0.37$, $+0.42$]\\
            & 30 & $-33.88$ & $+0.17$ & $-11.12$ & $-0.01$ & $+0.18$ [$+0.17$, $+0.20$]\\
\midrule
$8\times4$  & 0  & $-1.97$ & $+2.95$ & $+2.23$ & $+2.30$ & $+0.65$ [$+0.58$, $+0.72$]\\
            & 10 & $-2.57$ & $+4.59$ & $+4.39$ & $+4.54$ & $+0.05$ [$-0.03$, $+0.13$]\\
            & 20 & $-2.83$ & $+5.44$ & $+4.70$ & $+4.91$ & $+0.54$ [$+0.44$, $+0.63$]\\
            & 30 & $-2.96$ & $+5.85$ & $+4.61$ & $+4.75$ & $+1.11$ [$+0.99$, $+1.22$]\\
\bottomrule
\end{tabular}
\end{table}

\subsection{Discussion and limitations}
A learned estimator does not help across frequency within a 3-MHz block. There the channel is almost a
low-rank linear family, LMMSE with half a million training vectors is effectively optimal, and an MLP adds
approximation error and about 8--32 times the MACs. Its one advantage is robustness: without an SNR input
it stays within about 0.3\,dB of matched LMMSE up to 20\,dB, whereas LMMSE loses 2\,dB or more once the
SNR is assumed 10\,dB or more too low. A learned prior helps
where the set of channels is nonlinear and the pilots leave the channel underdetermined: across antennas,
most clearly with pilots on a subset of them.

\emph{Limitations.} All runs come from one room, and the test run covers the same area as the training
runs, so the intervals reflect noise and block sampling within one run, not a new environment. The pilot
noise is synthetic and white, and the measured CSI used as reference has its own error. The joint setting uses 16 of the 64 subcarriers and only 15\,404 training samples for
a 512-dimensional mixture: EM stopped after two iterations without changing the $k$-means partition, so
this GMM is a partition with one shrunk covariance per cell.

%=====================================================================================================
\section{Conclusion}
On measured indoor channels, a learned estimator did not help over LMMSE for reconstructing a
64-subcarrier block from 16, 8 or 4 noisy pilots: in a single held-out evaluation none of four MLP variants
was better in any of 84 settings, while LMMSE gained 1.7--9.7\,dB over linear interpolation. A learned prior did
help for joint reconstruction over antennas and subcarriers: a GMM estimator was better than joint LMMSE
in all eight settings, by up to 5.9\,dB with pilots on 8 of 32 antennas. For a receiver: within a narrow
block, use LMMSE with trained statistics (about 2.5\,dB loss per halving of the pilots); if the noise
level is unreliable, a blind MLP costs about 0.3\,dB or less up to 20\,dB; the mixture costs up to 64 times the
computation and does not yet make pilots on a quarter of the antennas accurate. Future work: a second
environment
(dichasus-015x), structured (Kronecker) mixture covariances, wider blocks, and bit-error rate.

\footnotesize
\bibliographystyle{IEEEtran}
\bibliography{reference}

\end{document}
